\documentclass[11pt]{article}
\usepackage[margin=1in]{geometry}
\usepackage{amsmath,amssymb,amsthm,booktabs,microtype}
\usepackage[hidelinks]{hyperref}
\newtheorem{theorem}{Theorem}[section]
\newtheorem{corollary}[theorem]{Corollary}
\newtheorem{remark}[theorem]{Remark}
\newcommand{\SU}{\mathrm{SU}(2)}
\title{Continuous Nonabelian Collar Charging\\
\large An exact geometric charge and an ultraviolet compatibility no-go}
\author{}
\date{20 July 2026}
\begin{document}
\maketitle

\begin{abstract}
We give a deterministic, continuous $\SU$ counterpart to center-contour
charging.  For plaquettes $P_i=US_i$ incident on a selected link, relative
staple distances equal relative plaquette distances for every value of the
selected link.  Failure of a common staple ball of radius $\eta$ therefore
forces an incident curvature above $\eta/2$ and total incident squared
curvature above $\eta^2/2$.  Combining this exact charge with the coherent-
collar requirement $\eta<\rho/\sqrt m$ proves that the strongest local charge
at the frozen ultraviolet witness has exponent only $0.9339$, versus the
required rate near seven.  Thus the simplest one-link charged-collar repair is
quantitatively incompatible with the declared parameters.
\end{abstract}

\section{Relative geometry}

Orient the $m$ plaquettes incident on a link so that
\[
 P_i=US_i,
\]
where $U$ is the selected-link variable and $S_i$ is its staple.  Let $d$ be
the bi-invariant geodesic distance on $\SU$.

\begin{theorem}[Link-uniform collar charge]
If the staple family is not contained in any common geodesic ball of radius
$\eta$, then for every $U\in\SU$,
\begin{align}
 \max_i d(I,P_i)&>\frac\eta2,\label{eq:single}\\
 \sum_{i=1}^m d(I,P_i)^2&>\frac{\eta^2}{2}.
 \label{eq:energy}
\end{align}
Both statements are gauge covariant.
\end{theorem}
\begin{proof}
Bi-invariance gives
\[
 d(P_i,P_j)=d(US_i,US_j)=d(S_i,S_j).
\]
If the staple diameter were at most $\eta$, choosing any staple as a center
would put the family inside an $\eta$-ball.  Collar failure therefore supplies
$i,j$ with $D=d(P_i,P_j)>\eta$.  Set
$a=d(I,P_i)$ and $b=d(I,P_j)$.  The triangle inequality gives $a+b\geq D$,
so $\max(a,b)>\eta/2$.  Moreover
\[
 a^2+b^2\geq\frac{(a+b)^2}{2}>\frac{\eta^2}{2},
\]
which proves \eqref{eq:energy}.  Gauge covariance follows because all
distances are bi-invariant.
\end{proof}

The theorem is stronger than a center-sector statement: no discreteness,
commutativity, or coordinate chart is used.

\section{Compatibility with the good collar}

The coherent conditional theorem requires
\begin{equation}
 \eta<\frac{\rho}{\sqrt m}.
 \label{eq:good}
\end{equation}
At the limiting admissible radius, \eqref{eq:single} charges only curvature
\[
 r_{\max}=\frac{\rho}{2\sqrt m},
\]
and \eqref{eq:energy}, inserted into a heat-kernel exponential
$e^{-E/(2t)}$, supplies exponent at most
\begin{equation}
 K_{\rm energy,max}=\frac{\rho^2}{4mt}.
 \label{eq:kmax}
\end{equation}

For a three-dimensional isotropic Gaussian radial event, put
\[
 J_3(s)=\begin{cases}
 \tfrac12(s-3-3\log(s/3)),&s>3,\\
 0,&s\leq3.
 \end{cases}
\]
The single-charge scaled radius is $s_{\max}=\rho^2/(4mt)$.

\begin{corollary}[Local-charging incompatibility]
At $t=3\times10^{-6}$, $\rho=0.0082$, and $m=6$,
\[
 s_{\max}=K_{\rm energy,max}=0.933888\ldots,
 \qquad J_3(s_{\max})=0.
\]
Consequently local collar charging cannot supply either the declared
$K=7$ or the exact Stage 14 minimum $K=6.665348\ldots$.
\end{corollary}

For a total-energy target $K$, compatibility requires
\begin{equation}
 t\leq \frac{\rho^2}{4mK}.
\end{equation}
At $K=7$ this is $t\leq4.00238\times10^{-7}$.  Solving
$J_3(s)=7$ gives $s=23.12719\ldots$ and the stricter radial requirement
$t\leq1.21142\times10^{-7}$.

\section{Interpretation and frontier}

The result does not say collar failure is unchargeable.  It proves the exact
opposite and then shows that this local charge is too weak at the chosen scale.
The charge must become extensive through multi-link integration, a coarser
collar, a typed center/topological contour, or a smaller-time witness.

No probability estimate for the complete interacting measure is inferred
from deterministic geometry alone.  No RG iteration, continuum limit, or
mass gap follows.  The result is a C2 restricted theorem and no-go
calibration.

\end{document}
